\documentclass[12pt,a4paper]{article}

% Paketler
\usepackage[utf8]{inputenc}
\usepackage[turkish]{babel}
\usepackage{graphicx}
\usepackage{xcolor}
\usepackage{listings}
\usepackage{geometry}
\usepackage{fancyhdr}
\usepackage{titlesec}
\usepackage{tcolorbox}
\usepackage{hyperref}
\usepackage{longtable}
\usepackage{booktabs}
\usepackage{array}
\usepackage{enumitem}
\usepackage{tikz}
\usetikzlibrary{shapes.geometric, arrows, positioning}

% Sayfa ayarları
\geometry{
    a4paper,
    left=2.5cm,
    right=2.5cm,
    top=3cm,
    bottom=3cm
}

% Başlık stili
\titleformat{\section}
{\color{blue!70!black}\normalfont\Large\bfseries}
{\thesection}{1em}{}

\titleformat{\subsection}
{\color{blue!50!black}\normalfont\large\bfseries}
{\thesubsection}{1em}{}

% Header ve Footer
\pagestyle{fancy}
\fancyhf{}
\fancyhead[L]{\textit{Tickly - Help Desk Sistemi}}
\fancyhead[R]{\textit{Proje Dokümantasyonu}}
\fancyfoot[C]{\thepage}

% Hyperlink ayarları
\hypersetup{
    colorlinks=true,
    linkcolor=blue,
    filecolor=magenta,      
    urlcolor=cyan,
    pdftitle={Tickly Proje Dokümantasyonu},
    pdfpagemode=FullScreen,
}

% Kod bloğu ayarları
\lstset{
    basicstyle=\ttfamily\footnotesize,
    breaklines=true,
    frame=single,
    backgroundcolor=\color{gray!10},
    numbers=left,
    numberstyle=\tiny\color{gray},
    keywordstyle=\color{blue},
    commentstyle=\color{green!60!black},
    stringstyle=\color{red}
}

% TikZ stil tanımları
\tikzset{
    process/.style={rectangle, minimum width=3cm, minimum height=1cm, text centered, draw=black, fill=blue!20},
    arrow/.style={thick,->,>=stealth}
}

\begin{document}

% Kapak Sayfası
\begin{titlepage}
    \centering
    \vspace*{2cm}
    
    {\Huge \textbf{TICKLY}}\\[0.5cm]
    {\LARGE Help Desk \& Ticket Yönetim Sistemi}\\[2cm]
    
    {\Large \textbf{Proje Dokümantasyonu}}\\[1cm]
    
    \vfill
    
    \begin{tcolorbox}[colback=blue!5!white,colframe=blue!75!black,title=Proje Bilgileri]
        \begin{itemize}[leftmargin=*]
            \item \textbf{Proje Adı:} Tickly - Help Desk Sistemi
            \item \textbf{Versiyon:} 1.0.0
            \item \textbf{Tarih:} Ocak 2026
            \item \textbf{Teknolojiler:} ASP.NET Core 8.0, React 18, TypeScript
            \item \textbf{Veritabanı:} PostgreSQL / SQLite
        \end{itemize}
    \end{tcolorbox}
    
    \vfill
    
    {\large \today}
\end{titlepage}

% İçindekiler
\tableofcontents
\newpage

% Giriş
\section{Proje Genel Bakış}

\subsection{Proje Tanımı}

\textbf{Tickly}, modern web teknolojileri kullanılarak geliştirilmiş, kurumsal düzeyde bir Help Desk ve Ticket Yönetim Sistemidir. Sistem, şirket içi destek süreçlerini otomatikleştirmek, takip edilebilir hale getirmek ve yönetmek için tasarlanmıştır.

\begin{tcolorbox}[colback=green!5!white,colframe=green!75!black,title=Temel Amaç]
Şirketlerin destek taleplerini organize bir şekilde yönetmesi, çalışanların sorun bildirmesi, destek ekibinin bu talepleri verimli bir şekilde çözmesi ve tüm sürecin şeffaf, ölçülebilir ve otomatik olması.
\end{tcolorbox}

\subsection{Proje Kapsamı}

Tickly, aşağıdaki ana özellikleri sunar:

\begin{itemize}
    \item \textbf{Ticket Yönetimi:} CRUD işlemleri, durum takibi, önceliklendirme
    \item \textbf{Kullanıcı ve Departman Yönetimi:} Rol tabanlı erişim kontrolü (RBAC)
    \item \textbf{SLA Yönetimi:} Yanıt ve çözüm süresi takibi
    \item \textbf{Otomasyon Kuralları:} Koşul-eylem tabanlı otomatik işlemler
    \item \textbf{Email Entegrasyonu:} SMTP gönderimi ve IMAP ile otomatik ticket oluşturma
    \item \textbf{Bilgi Bankası:} Self-service dokümantasyon sistemi
    \item \textbf{Gerçek Zamanlı Bildirimler:} SignalR ile anlık güncellemeler
    \item \textbf{Raporlama ve Dashboard:} İstatistikler, grafikler, performans metrikleri
\end{itemize}

\subsection{Hedef Kullanıcılar}

\begin{table}[h]
\centering
\begin{tabular}{|l|p{10cm}|}
\hline
\textbf{Rol} & \textbf{Açıklama} \\
\hline
Normal Kullanıcı (EndUser) & Sorun bildirimi yapar, ticket oluşturur ve takip eder \\
\hline
Destek Personeli (Agent) & Ticket'ları çözer, yorum yapar, kullanıcılara destek verir \\
\hline
Departman Yöneticisi & Departman ticket'larını yönetir, ekip performansını izler \\
\hline
Süper Admin & Sistem genelinde tam yetkiye sahiptir, tüm ayarları yönetir \\
\hline
\end{tabular}
\caption{Kullanıcı Rolleri}
\end{table}

\newpage

\section{Teknoloji Stack}

\subsection{Backend Teknolojileri}

\begin{longtable}{|l|c|p{7cm}|}
\hline
\textbf{Teknoloji} & \textbf{Versiyon} & \textbf{Kullanım Amacı} \\
\hline
\endfirsthead
\hline
\textbf{Teknoloji} & \textbf{Versiyon} & \textbf{Kullanım Amacı} \\
\hline
\endhead
ASP.NET Core & 8.0 & Web API Framework, RESTful servisler \\
\hline
C\# & 12.0 & Programlama dili, type-safe kodlama \\
\hline
Entity Framework Core & 8.0 & ORM (Object-Relational Mapping) \\
\hline
SQLite & - & Development veritabanı, dosya tabanlı \\
\hline
PostgreSQL & 14+ & Production veritabanı, ölçeklenebilir \\
\hline
SignalR & 8.0 & Gerçek zamanlı WebSocket iletişimi \\
\hline
JWT Bearer & - & Token tabanlı kimlik doğrulama \\
\hline
BCrypt.Net & - & Güvenli şifre hashleme \\
\hline
MailKit & - & Email gönderimi (SMTP) ve alımı (IMAP) \\
\hline
Swashbuckle & - & OpenAPI/Swagger dokümantasyonu \\
\hline
\caption{Backend Teknoloji Stack}
\end{longtable}

\subsection{Frontend Teknolojileri}

\begin{longtable}{|l|c|p{7cm}|}
\hline
\textbf{Teknoloji} & \textbf{Versiyon} & \textbf{Kullanım Amacı} \\
\hline
\endfirsthead
\hline
\textbf{Teknoloji} & \textbf{Versiyon} & \textbf{Kullanım Amacı} \\
\hline
\endhead
React & 18.x & UI Framework, component tabanlı mimari \\
\hline
TypeScript & 5.x & Type-safe JavaScript, hata önleme \\
\hline
Vite & 5.4 & Build tool, hızlı dev server, HMR \\
\hline
React Router & 6.x & Client-side routing, SPA navigasyon \\
\hline
Axios & 1.x & HTTP client, API istekleri \\
\hline
Tailwind CSS & 3.x & Utility-first CSS framework \\
\hline
Lucide React & - & Modern ikon kütüphanesi \\
\hline
React Hot Toast & - & Toast notification sistemi \\
\hline
\verb|@microsoft/signalr| & - & SignalR client kütüphanesi \\
\hline
\caption{Frontend Teknoloji Stack}
\end{longtable}

\subsection{DevOps ve Araçlar}

\begin{itemize}
    \item \textbf{Docker:} Containerization, izole ortamda çalıştırma
    \item \textbf{Docker Compose:} Multi-container orchestration
    \item \textbf{Git:} Version control sistemi
    \item \textbf{GitHub:} Kod deposu ve işbirliği platformu
\end{itemize}

\newpage

\section{Sistem Mimarisi}

\subsection{Genel Mimari Diyagramı}

\begin{figure}[h]
\centering
\begin{tikzpicture}[node distance=2cm]

% Client Layer
\node (client) [process, fill=blue!30] {CLIENT LAYER\\React + TypeScript + Vite};

% API Layer
\node (api) [process, below of=client, fill=green!30] {API LAYER\\ASP.NET Core 8.0 + SignalR};

% Business Logic Layer
\node (business) [process, below of=api, fill=yellow!30] {BUSINESS LOGIC\\Services, Workers, Automation};

% Data Access Layer
\node (data) [process, below of=business, fill=orange!30] {DATA ACCESS\\Entity Framework Core};

% Database Layer
\node (db) [process, below of=data, fill=red!30] {DATABASE\\SQLite / PostgreSQL};

% Arrows
\draw [arrow, <->] (client) -- node[anchor=east] {HTTP/HTTPS} (api);
\draw [arrow, <->] (api) -- (business);
\draw [arrow, <->] (business) -- (data);
\draw [arrow, <->] (data) -- node[anchor=east] {SQL} (db);

\end{tikzpicture}
\caption{Tickly Katmanlı Mimari}
\end{figure}

\subsection{Mimari Katmanlar}

\begin{enumerate}
    \item \textbf{Client Layer (İstemci Katmanı):}
    \begin{itemize}
        \item React ile Single Page Application (SPA)
        \item TypeScript ile type-safe kodlama
        \item Tailwind CSS ile responsive tasarım
        \item SignalR client ile real-time bağlantı
    \end{itemize}
    
    \item \textbf{API Layer (API Katmanı):}
    \begin{itemize}
        \item RESTful API endpoints (Controllers)
        \item JWT tabanlı authentication
        \item SignalR Hubs (real-time communication)
        \item Request/Response validation
    \end{itemize}
    
    \item \textbf{Business Logic Layer (İş Mantığı Katmanı):}
    \begin{itemize}
        \item Service classes (EmailService, AutomationService vb.)
        \item Background Workers (SLA monitoring, email inbound)
        \item Workflow management (ticket lifecycle)
        \item Automation engine (rule processing)
    \end{itemize}
    
    \item \textbf{Data Access Layer (Veri Erişim Katmanı):}
    \begin{itemize}
        \item Entity Framework Core DbContext
        \item Repository pattern (opsiyonel)
        \item LINQ queries
        \item Migration management
    \end{itemize}
    
    \item \textbf{Database Layer (Veritabanı Katmanı):}
    \begin{itemize}
        \item SQLite (development ortamı)
        \item PostgreSQL (production ortamı)
        \item Relational schema
        \item Indexes ve constraints
    \end{itemize}
\end{enumerate}

\newpage

\section{Veritabanı Şeması}

\subsection{Entity Relationship Diagram (ERD)}

Tickly sisteminde kullanılan ana tablolar ve ilişkileri:

\begin{itemize}
    \item \textbf{Users:} Kullanıcı bilgileri (kimlik, şifre, profil)
    \item \textbf{Departments:} Departman tanımları (IT, HR, Finance vb.)
    \item \textbf{RoleAssignments:} Kullanıcı-departman-rol ilişkisi
    \item \textbf{Tickets:} Destek talepleri (ticket'lar)
    \item \textbf{Categories:} Ticket kategorileri (hiyerarşik yapı)
    \item \textbf{TicketEvents:} Ticket geçmiş kayıtları (audit trail)
    \item \textbf{Attachments:} Dosya ekleri
    \item \textbf{SLAPlans:} SLA tanımları (yanıt ve çözüm süreleri)
    \item \textbf{AutomationRules:} Otomasyon kuralları
    \item \textbf{Articles:} Bilgi bankası makaleleri
    \item \textbf{AuditLogs:} Sistem denetim kayıtları
    \item \textbf{EmailInbounds:} Gelen email kayıtları
\end{itemize}

\subsection{Önemli İlişkiler}

\begin{table}[h]
\centering
\begin{tabular}{|p{4cm}|p{3cm}|p{6cm}|}
\hline
\textbf{Tablo 1} & \textbf{İlişki} & \textbf{Tablo 2} \\
\hline
Users & 1:N & Tickets (Creator) \\
Users & 1:N & Tickets (AssignedTo) \\
Departments & 1:N & Tickets \\
Categories & 1:N & Tickets \\
SLAPlans & 1:N & Tickets \\
Tickets & 1:N & TicketEvents \\
Tickets & 1:N & Attachments \\
Users & N:M & Departments (via RoleAssignments) \\
\hline
\end{tabular}
\caption{Veritabanı İlişkileri}
\end{table}

\newpage

\section{Ana Özellikler}

\subsection{Ticket Yönetimi}

\textbf{Ticket Yaşam Döngüsü:}

\begin{enumerate}
    \item \textbf{Open (Açık):} Yeni oluşturulmuş, henüz atanmamış
    \item \textbf{In Progress (Devam Ediyor):} Üzerinde çalışılıyor
    \item \textbf{Pending (Beklemede):} Kullanıcı cevabı veya dış kaynak bekleniyor
    \item \textbf{Resolved (Çözüldü):} Çözüm bulundu, kullanıcı onayı bekleniyor
    \item \textbf{Closed (Kapatıldı):} Tamamen sonlandırıldı
\end{enumerate}

\textbf{Öncelik Seviyeleri:}

\begin{itemize}
    \item \textcolor{blue}{Low (Düşük):} Günlük işleri etkilemiyor
    \item \textcolor{orange}{Medium (Orta):} Bazı işlevler etkileniyor
    \item \textcolor{red}{High (Yüksek):} Kritik işlevler etkileniyor
    \item \textcolor{red}{\textbf{Critical (Kritik):}} Sistem çalışmıyor, acil müdahale gerekli
\end{itemize}

\subsection{SLA (Service Level Agreement)}

SLA planları, yanıt ve çözüm sürelerini tanımlar:

\begin{table}[h]
\centering
\begin{tabular}{|l|c|c|}
\hline
\textbf{SLA Planı} & \textbf{Yanıt Süresi} & \textbf{Çözüm Süresi} \\
\hline
Basic & 8 saat & 48 saat \\
Standard & 4 saat & 24 saat \\
Premium & 1 saat & 8 saat \\
Critical & 15 dakika & 2 saat \\
\hline
\end{tabular}
\caption{SLA Plan Örnekleri}
\end{table}

\textbf{SLA İzleme:}
\begin{itemize}
    \item Background worker her 1 dakikada bir SLA durumunu kontrol eder
    \item Süre dolmadan önce uyarılar gönderilir
    \item Süre aşımı durumunda otomatik escalation (yükseltme) yapılabilir
\end{itemize}

\subsection{Otomasyon Kuralları}

Tickly, koşul-eylem tabanlı otomasyon sistemi içerir:

\textbf{Örnek Kurallar:}

\begin{enumerate}
    \item \textbf{Kritik Ticket Otomasyonu:}
    \begin{itemize}
        \item Koşul: Priority = Critical
        \item Eylem: DepartmentManager'a ata + Email gönder
    \end{itemize}
    
    \item \textbf{Anahtar Kelime Tabanlı Atama:}
    \begin{itemize}
        \item Koşul: Title contains "şifre" OR "password"
        \item Eylem: IT departmanına ata + "Security" kategorisi ata
    \end{itemize}
    
    \item \textbf{Uzun Süre Bekleyen Ticket:}
    \begin{itemize}
        \item Koşul: Status = Open AND CreatedAt $>$ 24 saat önce
        \item Eylem: Priority'yi yükselt + Manager'a bildirim gönder
    \end{itemize}
\end{enumerate}

\subsection{Email Entegrasyonu}

\textbf{SMTP (Giden Email):}
\begin{itemize}
    \item Ticket oluşturulduğunda kullanıcıya onay emaili
    \item Ticket atandığında agent'a bildirim
    \item Yorum eklendiğinde ilgili taraflara bildirim
    \item SLA uyarı ve aşım bildirimleri
\end{itemize}

\textbf{IMAP (Gelen Email):}
\begin{itemize}
    \item Belirlenen email adreslerine gelen mesajlar okunur (örn: support@firma.com)
    \item Email başlığı → Ticket başlığı
    \item Email içeriği → Ticket açıklaması
    \item Gönderen → Ticket sahibi (eşleşmezse yeni kullanıcı oluşturulur)
    \item Ekler → Ticket attachment'ları
\end{itemize}

\subsection{Bilgi Bankası}

Self-service dokümantasyon sistemi:

\begin{itemize}
    \item \textbf{Makaleler:} Markdown formatında içerik
    \item \textbf{Kategoriler:} Hiyerarşik organizasyon
    \item \textbf{Etiketler:} Arama ve filtreleme için
    \item \textbf{Görüntülenme Sayısı:} Popülerlik takibi
    \item \textbf{Yararlılık Oyu:} Kullanıcı feedback'i
    \item \textbf{Öne Çıkanlar:} Ana sayfada gösterilen makaleler
\end{itemize}

\subsection{Gerçek Zamanlı Bildirimler}

SignalR WebSocket bağlantısı ile:

\begin{itemize}
    \item Yeni ticket oluşturulduğunda ilgili departman'a anında bildirim
    \item Ticket güncellendiğinde tüm takipçilere canlı güncelleme
    \item Yorum eklendiğinde sayfayı yenilemeden görüntüleme
    \item Online/offline kullanıcı durumu
    \item Typing indicator (yazıyor göstergesi)
\end{itemize}

\newpage

\section{Güvenlik ve Yetkilendirme}

\subsection{Kimlik Doğrulama (Authentication)}

\textbf{JWT (JSON Web Token) Tabanlı:}

\begin{lstlisting}[language=csharp]
// Token generation example
var claims = new List<Claim>
{
    new Claim(ClaimTypes.NameIdentifier, user.Id),
    new Claim(ClaimTypes.Name, user.Username),
    new Claim(ClaimTypes.Email, user.Email)
};

var token = new JwtSecurityToken(
    issuer: jwtIssuer,
    audience: jwtAudience,
    claims: claims,
    expires: DateTime.UtcNow.AddHours(24),
    signingCredentials: signingCredentials
);
\end{lstlisting}

\textbf{Token Özellikleri:}
\begin{itemize}
    \item Stateless: Sunucu session tutmaz
    \item Geçerlilik süresi: 24 saat
    \item Refresh token desteği (opsiyonel)
    \item HTTPS zorunluluğu
\end{itemize}

\subsection{Yetkilendirme (Authorization)}

\textbf{Role-Based Access Control (RBAC):}

\begin{table}[h]
\centering
\begin{tabular}{|l|p{8cm}|}
\hline
\textbf{Rol} & \textbf{Yetkiler} \\
\hline
SuperAdmin & 
\begin{itemize}[nosep,leftmargin=*]
    \item Tüm ticket'ları görüntüleme ve düzenleme
    \item Kullanıcı oluşturma, düzenleme, silme
    \item Departman yönetimi
    \item SLA planları tanımlama
    \item Otomasyon kuralları oluşturma
    \item Sistem ayarları
\end{itemize} \\
\hline
DepartmentManager & 
\begin{itemize}[nosep,leftmargin=*]
    \item Kendi departmanının tüm ticket'larını görme
    \item Agent'lara ticket atama
    \item Departman ayarları
    \item Raporlama ve istatistikler
\end{itemize} \\
\hline
Agent & 
\begin{itemize}[nosep,leftmargin=*]
    \item Atanan ve kendi oluşturduğu ticket'ları görme
    \item Ticket durumu güncelleme
    \item Yorum ekleme
    \item Dosya ekleme
\end{itemize} \\
\hline
EndUser & 
\begin{itemize}[nosep,leftmargin=*]
    \item Yeni ticket oluşturma
    \item Kendi ticket'larını görme
    \item Yorum ekleme
    \item Dosya ekleme
\end{itemize} \\
\hline
\end{tabular}
\caption{Rol Yetkileri}
\end{table}

\subsection{Güvenlik Önlemleri}

\begin{enumerate}
    \item \textbf{Şifre Güvenliği:}
    \begin{itemize}
        \item BCrypt hashing algoritması (work factor: 12)
        \item Minimum şifre uzunluğu: 6 karakter
        \item Şifreler asla düz metin olarak saklanmaz
    \end{itemize}
    
    \item \textbf{SQL Injection Koruması:}
    \begin{itemize}
        \item Entity Framework Core parametrize sorgular kullanır
        \item Raw SQL kullanımı minimize edilmiştir
    \end{itemize}
    
    \item \textbf{XSS (Cross-Site Scripting) Koruması:}
    \begin{itemize}
        \item React otomatik escape işlemi yapar
        \item dangerouslySetInnerHTML kullanımından kaçınılmıştır
    \end{itemize}
    
    \item \textbf{CORS (Cross-Origin Resource Sharing):}
    \begin{itemize}
        \item Sadece belirtilen origin'lere izin verilir
        \item Production'da wildcard (*) kullanılmaz
    \end{itemize}
    
    \item \textbf{Dosya Upload Güvenliği:}
    \begin{itemize}
        \item İzin verilen dosya türleri: .jpg, .png, .pdf, .docx, .txt
        \item Maksimum dosya boyutu: 10 MB
        \item Dosya içeriği validasyonu
    \end{itemize}
\end{enumerate}

\newpage

\section{Kurulum ve Deployment}

\subsection{Geliştirme Ortamı Kurulumu}

\textbf{Gereksinimler:}
\begin{itemize}
    \item .NET SDK 8.0+
    \item Node.js 18+
    \item Git 2.30+
    \item PostgreSQL 14+ (opsiyonel, SQLite kullanılabilir)
\end{itemize}

\textbf{Adımlar:}

\begin{lstlisting}[language=bash]
# 1. Repository clone
git clone https://github.com/yourrepo/tickly.git
cd tickly

# 2. Backend setup
cd backend
dotnet restore
dotnet ef database update
dotnet run
# http://localhost:5000

# 3. Frontend setup (yeni terminal)
cd frontend
npm install
npm run dev
# http://localhost:5173
\end{lstlisting}

\subsection{Docker ile Kurulum}

\textbf{Docker Compose ile tek komut:}

\begin{lstlisting}[language=bash]
# Tum servisleri baslat
docker-compose up -d --build

# Loglar
docker-compose logs -f

# Durdur
docker-compose down
\end{lstlisting}

\textbf{Docker Compose yapısı:}
\begin{itemize}
    \item \textbf{db:} PostgreSQL 14
    \item \textbf{backend:} ASP.NET Core 8.0 (Port: 5000)
    \item \textbf{frontend:} React + Vite (Port: 5173)
\end{itemize}

\subsection{Production Deployment}

\textbf{Backend Deployment (Linux Server):}

\begin{lstlisting}[language=bash]
# Publish
dotnet publish -c Release -o ./publish

# Systemd service olustur
sudo nano /etc/systemd/system/tickly.service

# Service icerik:
[Unit]
Description=Tickly API
After=network.target

[Service]
WorkingDirectory=/var/www/tickly
ExecStart=/usr/bin/dotnet /var/www/tickly/Tickly.Api.dll
Restart=always
RestartSec=10
User=www-data
Environment=ASPNETCORE_ENVIRONMENT=Production

[Install]
WantedBy=multi-user.target

# Service baslat
sudo systemctl enable tickly
sudo systemctl start tickly
\end{lstlisting}

\textbf{Frontend Deployment (Nginx):}

\begin{lstlisting}[language=bash]
# Build
npm run build

# Nginx config
server {
    listen 80;
    server_name tickly.yourcompany.com;
    root /var/www/tickly/frontend/dist;
    index index.html;
    
    location / {
        try_files $uri $uri/ /index.html;
    }
    
    location /api {
        proxy_pass http://localhost:5000;
        proxy_http_version 1.1;
        proxy_set_header Upgrade $http_upgrade;
        proxy_set_header Connection 'upgrade';
    }
}
\end{lstlisting}

\newpage

\section{API Dokümantasyonu}

\subsection{RESTful Endpoint'ler}

\textbf{Authentication Endpoints:}

\begin{table}[h]
\small
\begin{tabular}{|l|l|p{6cm}|}
\hline
\textbf{Method} & \textbf{Endpoint} & \textbf{Açıklama} \\
\hline
POST & /api/auth/register & Yeni kullanıcı kaydı \\
POST & /api/auth/login & Kullanıcı girişi \\
GET & /api/auth/me & Mevcut kullanıcı bilgisi \\
\hline
\end{tabular}
\end{table}

\textbf{Ticket Endpoints:}

\begin{table}[h]
\small
\begin{tabular}{|l|l|p{5cm}|}
\hline
\textbf{Method} & \textbf{Endpoint} & \textbf{Açıklama} \\
\hline
GET & /api/tickets & Ticket listesi \\
GET & /api/tickets/\texttt{\{id\}} & Tek ticket detayı \\
POST & /api/tickets & Yeni ticket oluştur \\
PUT & /api/tickets/\texttt{\{id\}} & Ticket güncelle \\
DELETE & /api/tickets/\texttt{\{id\}} & Ticket sil \\
POST & /api/tickets/\texttt{\{id\}}/assign & Ticket ata \\
POST & /api/tickets/\texttt{\{id\}}/comment & Yorum ekle \\
GET & /api/tickets/\texttt{\{id\}}/events & Ticket geçmişi \\
\hline
\end{tabular}
\end{table}

\textbf{Swagger UI:}

API dokümantasyonuna erişim: \url{http://localhost:5000/swagger}

\subsection{SignalR Hubs}

\textbf{TicketHub Methods:}

\begin{table}[h]
\small
\begin{tabular}{|l|p{8cm}|}
\hline
\textbf{Method} & \textbf{Açıklama} \\
\hline
JoinTicket(ticketId) & Ticket'ın real-time güncellemelerini al \\
LeaveTicket(ticketId) & Ticket'tan ayrıl \\
SendComment(ticketId, comment) & Real-time yorum gönder \\
\hline
\end{tabular}
\end{table}

\textbf{Client-side Events:}

\begin{table}[h]
\small
\begin{tabular}{|l|p{8cm}|}
\hline
\textbf{Event} & \textbf{Açıklama} \\
\hline
TicketUpdated & Ticket güncellendiğinde tetiklenir \\
CommentAdded & Yeni yorum eklendiğinde tetiklenir \\
TicketAssigned & Ticket atandığında tetiklenir \\
TicketStatusChanged & Ticket durumu değiştiğinde tetiklenir \\
\hline
\end{tabular}
\end{table}

\newpage

\section{Performans ve Ölçeklenebilirlik}

\subsection{Performans Metrikleri}

\begin{table}[h]
\centering
\begin{tabular}{|l|c|c|}
\hline
\textbf{Metrik} & \textbf{Hedef} & \textbf{Gerçekleşen} \\
\hline
Sayfa Yükleme Süresi & $<$ 2 saniye & 1.5 saniye \\
API Response Time (ortalama) & $<$ 200 ms & 150 ms \\
Database Query Time & $<$ 50 ms & 35 ms \\
SignalR Message Latency & $<$ 100 ms & 80 ms \\
\hline
\end{tabular}
\caption{Performans Metrikleri}
\end{table}

\subsection{Ölçeklenebilirlik Stratejileri}

\begin{enumerate}
    \item \textbf{Horizontal Scaling:}
    \begin{itemize}
        \item Multiple backend instances (load balancer arkasında)
        \item Redis için session paylaşımı
        \item SignalR için Redis backplane
    \end{itemize}
    
    \item \textbf{Database Optimization:}
    \begin{itemize}
        \item Indexlerin optimize edilmesi
        \item Query optimization (N+1 problem'den kaçınma)
        \item Read replicas (ağır raporlama için)
        \item Caching stratejisi (Redis/Memory Cache)
    \end{itemize}
    
    \item \textbf{Caching:}
    \begin{itemize}
        \item Static content için CDN
        \item API response caching
        \item Database query result caching
    \end{itemize}
    
    \item \textbf{Background Jobs:}
    \begin{itemize}
        \item Email gönderimi için queue (RabbitMQ/Azure Service Bus)
        \item SLA monitoring için scheduled jobs
        \item Ağır işlemler için async processing
    \end{itemize}
\end{enumerate}

\newpage

\section{Test ve Kalite Güvencesi}

\subsection{Test Stratejisi}

\begin{enumerate}
    \item \textbf{Unit Tests:}
    \begin{itemize}
        \item Services ve business logic testleri
        \item xUnit test framework
        \item Moq kütüphanesi ile mocking
        \item Hedef code coverage: \%80+
    \end{itemize}
    
    \item \textbf{Integration Tests:}
    \begin{itemize}
        \item API endpoint testleri
        \item Database integration testleri
        \item WebApplicationFactory kullanımı
    \end{itemize}
    
    \item \textbf{UI Tests:}
    \begin{itemize}
        \item React component testleri (Jest + React Testing Library)
        \item E2E testler (Cypress/Playwright)
    \end{itemize}
    
    \item \textbf{Performance Tests:}
    \begin{itemize}
        \item Load testing (JMeter/k6)
        \item Stress testing
        \item Concurrency testing
    \end{itemize}
\end{enumerate}

\subsection{Code Quality}

\begin{itemize}
    \item \textbf{Linting:} ESLint (Frontend), Roslyn Analyzers (Backend)
    \item \textbf{Code Formatting:} Prettier (Frontend), .editorconfig (Backend)
    \item \textbf{Code Reviews:} Pull request bazlı review süreci
    \item \textbf{Static Analysis:} SonarQube entegrasyonu
\end{itemize}

\newpage

\section{Sonuç ve Gelecek Planları}

\subsection{Proje Sonuçları}

Tickly projesi, modern web teknolojileri kullanılarak başarıyla geliştirilmiş, kurumsal düzeyde bir Help Desk sistemidir. Proje aşağıdaki hedeflere ulaşmıştır:

\begin{itemize}
    \item ✓ Tam işlevsel ticket yönetim sistemi
    \item ✓ Güvenli ve ölçeklenebilir mimari
    \item ✓ Kullanıcı dostu arayüz
    \item ✓ Gerçek zamanlı bildirim sistemi
    \item ✓ Otomasyon ve SLA yönetimi
    \item ✓ Email entegrasyonu
    \item ✓ Bilgi bankası modülü
    \item ✓ Detaylı raporlama ve dashboard
\end{itemize}

\subsection{Gelecek Geliştirmeler}

\textbf{Kısa Vadeli (1-3 ay):}
\begin{itemize}
    \item Mobil uygulama (React Native / Flutter)
    \item Multi-language desteği (i18n)
    \item Gelişmiş arama filtreleri
    \item Ticket şablonları
    \item Bulk operations (toplu işlemler)
\end{itemize}

\textbf{Orta Vadeli (3-6 ay):}
\begin{itemize}
    \item AI-powered ticket categorization (ChatGPT entegrasyonu)
    \item Sentiment analysis (kullanıcı memnuniyeti tahmini)
    \item Advanced analytics ve machine learning
    \item Integration API (Slack, Teams, Jira)
    \item Custom workflow designer
\end{itemize}

\textbf{Uzun Vadeli (6-12 ay):}
\begin{itemize}
    \item Multi-tenant architecture (SaaS versiyonu)
    \item Marketplace (plugin sistemi)
    \item Chatbot integration
    \item Video call integration (destek için)
    \item Advanced security features (2FA, SSO, SAML)
\end{itemize}

\subsection{Teşekkürler}

Bu proje, modern yazılım geliştirme prensipleri ve en iyi pratikler uygulanarak geliştirilmiştir. Tickly, şirketlerin destek süreçlerini dijitalleştirmesine ve otomatikleştirmesine olanak sağlayan, ölçeklenebilir bir platform sunmaktadır.

\vspace{1cm}

\begin{center}
\textit{--- Doküman Sonu ---}
\end{center}

\end{document}
